後一部未列入每週指定範圍，供想讀同一題目其他作者的人自取；老師稱這兩部選集為本領域廣泛實用的
參考書，「為了上這個課我還蠻辛苦的，準備了好幾個禮拜找這些文獻然後把它排一排」。當天另補一篇
中文文章進參考資料，即陳文政、陳偉杰、莊旻達、王上維四人合著的〈美國司法違憲審查正當性論辯脈絡
之分析〉。\footcite{chen2011}

第 4 週的 Bickel 用 GAF 219 至 225 頁的節錄，不用 \emph{The Least Dangerous Branch} 全文；\footcite{bickel2004}
老師說第 9 週的 Kramer 同此，取的是選集所收的節錄。\footnote{GAF 所收的 Kramer 為 2001 年〈We the Court〉
的節錄，見 \cite{kramer2001gaf}；課綱第 9 週所列為 Kramer 2012 年〈Judicial Supremacy and the End of
Judicial Restraint〉（100 Calif.\ L.\ Rev.\ 621），與 GAF 所收者不是同一篇。}老師自任前三次導讀（9 月 21 日、10 月 5 日、10 月 12 日），
10 月 19 日的 Waldron 那一場起可由同學認領，分配由隨班助教與班上協調，第 2 週再定。10 月 12 日
那一場要不要讀 Kelsen 與 Schmitt 的全譯還在考慮，Vinx 譯本兩篇合計約一百頁。\footcites[22--78]{kelsen2015}[79--124]{schmitt2015}
行政須依法律執行，老師以美國總統為對照：美國憲法要求總統忠實執行法律，總統雖與國會對立，仍未
公然違反國會通過的法律；國會以法律訂定的地名，總統也只能照用，未經法律列名的才有更改的餘地。
他由此批評 113 年憲判字第 9 號以「憲法未授權即不得立法」為違憲理由。該號判決理由認總統之憲法上
地位及職權屬憲法保留事項，「除憲法明文規定或授權外，尚非可由立法院制定法律，課予總統憲法所無
之義務」，立法院於未有憲法依據下以法律規定總統應為或不為一定行為者，逾越其憲法上職權，違反
權力分立原則。\footcite[理由第 81 段；另見第 91、116 段]{ccj113-9} 老師反問，美國國會以法律訂定
地名，美國憲法亦無授權之明文，依同一理由豈不也不得立法。他接著說明，該號判決以前法律人所理解的法律保留，是用來防止命令
侵奪法律的地位：一為法律優位，命令與法律牴觸者無效；二為保留給法律的事項，不得以命令規定；
三為法律得規定哪些事項，原則上由立法者依重要事項保留自行判斷，立法者認為重要而予以規定者，
命令即不得與之牴觸。
